% Hémi-cylindre
\begin{minipage}{0.72\linewidth}
    Un mince faisceau de lumière monochromatique bleue vient frapper un
    hémi-cylindre de Plexiglas comme le montre la figure ci-contre.
    \begin{enumerate}
        \item Calculer la valeur de l'indice du Plexiglas pour cette radiation
              monochromatique sachant que les angles d'incidence et de
              réfraction valent~:
              \[
                  i=\ang{65.0}
                  \qquad\text{et}\qquad
                  r=\ang{37.2}.
              \]
    \end{enumerate}
\end{minipage}
\hfill
\begin{minipage}{0.24\linewidth}
    \flushright
    \begin{tikzpicture}
        \def\r{2}
        \def\i{65}
        \draw[fill=non-photoblue,thick] (-\r,0) arc (180:360:\r) -- cycle;
        \draw[dashed] (0,-1.1*\r) -- (0,0.8*\r);
        \draw[blue,very thick,postaction={decorate},
              decoration={markings,mark=at position 0.5 with
              {\arrowreversed{Straight Barb}}}] (0,0) -- ++(90+\i:\r);
        \draw (90:0.4*\r) arc (90:90+\i:0.4*\r) node[above,pos=0.5] {$i$};
        \node[above left] at (\r,0) {air};
    \end{tikzpicture}
\end{minipage}

\begin{enumerate}
    \setcounter{enumi}{1}
    \item Calculer la valeur de l'angle de réfraction pour une radiation
          lumineuse jaune qui vient frapper l'hémi-cylindre dans les mêmes
          conditions que la radiation bleue ($i=\ang{65.0}$) sachant que pour
          la lumière jaune, l'indice du Plexiglas est égal à $1,48$.
    \item Tracer la marche des rayons bleu et jaune dans l'hémi-cylindre et
          conclure.
    \item L'hémi-cylindre permet-il de réaliser la dispersion d'une lumière
          blanche~?
\end{enumerate}
